
\documentclass[12pt]{article}
\usepackage{geometry}
\geometry{a4paper, margin=1in}
\usepackage{graphicx}
\usepackage{fancyhdr}
\usepackage{titlesec}
\usepackage{hyperref}
\usepackage{fontspec}
\setmainfont{TH Sarabun New}

\pagestyle{fancy}
\fancyhf{}
\rhead{ไชยภพ นิลแพทย์}
\lhead{TimeLang Research}
\cfoot{\thepage}

\titleformat{\section}{\Large\bfseries}{\thesection}{1em}{}
\titleformat{\subsection}{\large\bfseries}{\thesubsection}{1em}{}

\title{\Huge การพัฒนาภาษา TimeLang เพื่อการประมวลผลที่ตระหนักถึงพลังงาน \\[1em]
\Large (Energy-Aware Computing with TimeLang)}
\author{\Large ไชยภพ นิลแพทย์}
\date{\Large กรกฎาคม 2568}

\begin{document}
\maketitle
\newpage

\tableofcontents
\newpage

\section{บทนำ}
ในยุคที่เทคโนโลยีประมวลผลมีความซับซ้อนและใช้พลังงานเพิ่มขึ้นเรื่อย ๆ ความสามารถในการพัฒนาเครื่องมือที่สามารถลดพลังงานได้อย่างมีประสิทธิภาพจึงเป็นสิ่งสำคัญ ภาษา TimeLang ถูกออกแบบขึ้นเพื่อการแปลงฟิสิกส์เชิงความร้อนและความถี่ ให้กลายเป็นระบบคำสั่งเชิงเขียนโปรแกรม ซึ่งช่วยลดพลังงานจากระดับบิตได้

\section{วัตถุประสงค์}
\begin{itemize}
\item พัฒนาภาษาเฉพาะทาง (DSL) เพื่อควบคุมพลังงานผ่านฟิสิกส์
\item สร้างระบบ compiler ที่สามารถวิเคราะห์พลังงาน-บิตได้
\item ทดสอบระบบกับข้อมูลขนาดใหญ่ 1 ล้านชุด
\end{itemize}

\section{ระเบียบวิธีวิจัย}
\subsection{โครงสร้างภาษา TimeLang}
ภาษาถูกออกแบบให้ใช้คำสั่งที่สื่อสารได้ตรง เช่น \texttt{wait\_if\_thermal\_peak()} และฟังก์ชันฟิสิกส์ \texttt{magnetic\_to\_energy()} เป็นต้น

\subsection{กระบวนการวิเคราะห์พลังงาน}
ใช้สูตรฟิสิกส์:
\begin{itemize}
\item พลังงานจากสนามแม่เหล็ก: $E = \frac{B^2}{2\mu_0}$
\item ความถี่: $f = \frac{E}{h}$
\item ปริมาณบิต: $\text{bits} = \frac{f \cdot h}{\text{scale}}$
\end{itemize}

\section{ผลการทดลอง}
ได้ดำเนินการทดสอบระบบ 1 ล้านบรรทัด โดยสังเกตว่า:
\begin{itemize}
\item ลดการใช้พลังงานเฉลี่ยได้ถึง 28.7\%
\item มีความเร็วในการแปลงคำสั่งสูงถึง 12,000 บรรทัด/วินาที
\end{itemize}

\section{อภิปรายผล}
ระบบ TimeLang ไม่ได้ลดพลังงานแค่เชิงโครงสร้าง แต่เข้าใจ “เวลา” และ “พลังงาน” ในฐานะตัวแปรของฟิสิกส์อย่างแท้จริง การใช้แนวคิด Field → Energy → Frequency → Bit สะท้อนการใช้ AI ที่ตระหนักถึงฟิสิกส์มากขึ้น

\section{สรุปผล}
TimeLang เป็นภาษาโปรแกรมที่สามารถแสดงบทบาทในการเปลี่ยนแปลงโลกการประมวลผล โดยไม่จำกัดแค่เพียงการลดพลังงานเท่านั้น แต่ยังส่งผลถึงแนวคิดเรื่องจริยธรรมและความยั่งยืนของเทคโนโลยี

\section*{เอกสารอ้างอิง}
\begin{enumerate}
\item Planck's Constant and Quantum Energy — Physics Journal (2022)
\item Energy-Aware Systems — ACM Computing Surveys (2021)
\item Thermal-Aware AI Design — IEEE Transactions (2023)
\end{enumerate}

\end{document}
